% Final report -- IEEE 2-column conference format, max 6 pages (course requirement).
% Compile on Overleaf (IEEEtran is included there) or with any TeX distribution that has IEEEtran.
\documentclass[conference]{IEEEtran}
\usepackage{cite}
\usepackage{amsmath,amssymb}
\usepackage{graphicx}
\usepackage{booktabs}
\usepackage{xcolor}
\newcommand{\todo}[1]{\textcolor{red}{[TODO: #1]}}

\begin{document}

\title{Supervised, Self-Supervised and Language-Supervised Backbones for Domain Generalization on PACS}

\author{
\IEEEauthorblockN{M\'arton Tam\'as Magyar (TG2UTM), Daniel Brody (IASULA), Marco Buonfantino (FMDKOH)}
\IEEEauthorblockA{Budapest University of Technology and Economics -- Deep Learning, Project 5}
}

\maketitle

\begin{abstract}
\todo{4--6 sentences: problem (classifiers trained on photos fail on paintings, cartoons, sketches), what we did (reproduce ERM on PACS with DomainBed; compare ResNet-50, DINOv2 and CLIP backbones adapted by linear probing, LoRA and full fine-tuning; research question of Milestone 3), main numbers, main conclusion.}
\end{abstract}

\section{Introduction}
\todo{Domain generalization, PACS (7 classes, 4 domains), leave-one-domain-out protocol, why the pretraining of the backbone matters, contributions as a bullet list (one per milestone).}

\section{Related Work}
\todo{DomainBed and fair model selection~\cite{gulrajani2021domainbed}; PACS~\cite{li2017pacs}; DINOv2~\cite{oquab2023dinov2}; CLIP~\cite{radford2021clip}; LoRA~\cite{hu2022lora}; WiSE-FT~\cite{wortsman2022wiseft}.}

\section{Method}
\subsection{Protocol}
All experiments use DomainBed with ERM on PACS, leave-one-domain-out (4 test domains), a 20\% hold-out of every domain, and model selection on the training-domain validation set; the test domain is never used for selection or tuning. Results are the test-domain accuracy of the selected checkpoint. Milestone 2 uses one seed (trial seed 0, the same split as Milestone 1) for every method. \todo{Milestone 3: main experiment with 3 seeds, mean $\pm$ std.}
\subsection{Backbones and adaptation}
\todo{ResNet-50 (ImageNet, supervised), DINOv2 ViT-B/14 (self-supervised; features = 4 CLS tokens + mean patch token), CLIP ViT-B/16 (language-supervised, LAION-2B, input renormalization). Linear probe; LoRA rank 8 on the attention projections (identical 442k parameters on both ViTs; parametrization for CLIP's MultiheadAttention); full fine-tuning with 10$\times$ head learning rate. Table of hyperparameters and steps.}

\section{Milestone 1: Reproduction}
\begin{table}[t]
\centering
\caption{ERM with ResNet-50 on PACS, training-domain validation (accuracy \%).}
\label{tab:m1}
\begin{tabular}{lccccc}
\toprule
 & A & C & P & S & Avg\\
\midrule
DomainBed paper~\cite{gulrajani2021domainbed} & 84.7 & 80.8 & 97.2 & 79.3 & 85.5\\
Ours (seed 0) & 82.5 & 80.9 & 97.8 & 78.3 & 84.9\\
\bottomrule
\end{tabular}
\end{table}
Our single run is within 0.6 points of the paper on average and within 2.2 points on every domain (Table~\ref{tab:m1}), inside the 1--3 point range expected from random initialization, library versions and the absence of the paper's hyperparameter search (we use the default hyperparameters). \todo{Discuss the per-domain differences (Art, Sketch).}

\section{Milestone 2: Improving the Backbone}
\begin{table*}[t]
\centering
\caption{Test-domain accuracy (\%, seed 0) and cost. \todo{fill from M2\_results.ipynb}}
\label{tab:m2}
\begin{tabular}{llccccccccc}
\toprule
Backbone & Adaptation & A & C & P & S & Avg & ECE & Trainable params & VRAM (GB) & s/step\\
\midrule
ResNet-50 (M1) & full (ERM) & 82.5 & 80.9 & 97.8 & 78.3 & 84.9 & & & & \\
DINOv2 & linear probe & & & & & & & & & \\
DINOv2 & LoRA & & & & & & & & & \\
DINOv2 & full FT & & & & & & & & & \\
CLIP & zero-shot & & & & & & & 0 & & \\
CLIP & linear probe & & & & & & & & & \\
CLIP & LoRA & & & & & & & & & \\
CLIP & full FT & & & & & & & & & \\
\bottomrule
\end{tabular}
\end{table*}
\todo{Effect of the backbone (one factor at a time), effect of the adaptation method, accuracy vs.\ cost, which differences are large enough to be meaningful with one seed, computational budget of each method (steps, time).}

\section{Milestone 3: Research Question}
\todo{State the research question(s); WiSE-FT curve with $\alpha$ selected on training-domain validation; few-shot results; calibration (ECE, reliability). Main experiment with 3 seeds, mean $\pm$ std.}

\section{Error Analysis and Limitations}
\todo{Per-class accuracy and most frequent confusions per held-out domain (from the confusion matrices), example failures (e.g.\ sketches), over-confidence on the unseen domain; limitations: single dataset, no hyperparameter search, fewer steps for the ViTs, compute constraints.}

\section{Conclusion}
\todo{3--4 sentences.}

\bibliographystyle{IEEEtran}
\bibliography{refs}

\end{document}
